\section{Background and GPU Design}
\label{sec:design}

A symbolic state pairs a discrete location vector with a zone stored as a
$(m{+}1)\times(m{+}1)$ Difference Bound Matrix in canonical form. The successor pipeline for
one $(\text{state},\text{edge})$ pair is shown in Fig.~\ref{fig:kernel}: integer guard, clock
guard, resets, combined invariant, delay, invariant again, inactive-clock reset,
extrapolation, emptiness. Integer guards are evaluated \emph{host-side} while the work list
is built, so the kernel receives only pairs whose discrete part is already known feasible ---
a fact that matters in Sect.~\ref{sec:why}.

\begin{figure}[t]
\centering
\begin{tikzpicture}[font=\small,
  st/.style={draw, rounded corners=2pt, minimum height=7mm, minimum width=15mm, align=center,
             font=\scriptsize, fill=cA!10, draw=cA},
  host/.style={st, fill=cD!12, draw=cD},
  ar/.style={-{Latex[length=1.8mm]}, thick},
  ex/.style={-{Latex[length=1.8mm]}, cB, thick}]

\node[host] (ig) at (0,0) {int guard\\\textit{host-side}};
\node[st, right=4mm of ig] (cg) {clock\\guard};
\node[st, right=4mm of cg] (rs) {resets};
\node[st, right=4mm of rs] (iv) {invariant};
\node[st, right=4mm of iv] (dl) {delay};
\node[st, right=4mm of dl] (iv2) {invariant};
\node[st, right=4mm of iv2] (ri) {reset\\inactive};
\node[st, right=4mm of ri] (ex) {extrap.\\+ empty?};

\foreach \a/\b in {ig/cg, cg/rs, rs/iv, iv/dl, dl/iv2, iv2/ri, ri/ex} \draw[ar] (\a) -- (\b);
\foreach \n in {cg, iv, iv2, ex} \draw[ex] (\n.south) -- ++(0,-5mm);
\node[font=\scriptsize, cB] at (5.0,-0.95) {early exits (all require the DBM)};
\end{tikzpicture}
\caption{Per-work-item successor pipeline. Only the first stage is DBM-free, and it already
runs on the host, so a cheap GPU pre-filter has nothing left to reject.}
\label{fig:kernel}
\end{figure}

\subsection{Three DBM representations}
\label{sec:reps}

We implemented three mappings of a zone onto the GPU (Fig.~\ref{fig:reps}), all producing
identical, \uppaal{}-exact state counts.

\begin{description}
\item[R1: private dense matrix (baseline).] One work item per thread; the zone lives in a
per-thread array \texttt{int z[MAXDIM*MAXDIM]}. At $\mathit{MAXDIM}=24$ this is 2304\,B.
\item[R2: threadgroup-resident dense matrix.] One work item per thread, but the zone lives in
threadgroup (on-chip) memory, sized by the model's real $\mathit{dim}$ rather than
$\mathit{MAXDIM}$, with an odd stride to avoid bank conflicts.
\item[R3: row-distributed.] One work item per \emph{SIMD group}; lane $i$ owns row $i$ of the
matrix in registers. Per-lane footprint falls from $\mathit{dim}^2$ to $\mathit{dim}$ --- a
$24\times$ reduction. Floyd--Warshall obtains $d[k][\cdot]$ by \texttt{simd\_shuffle},
register-to-register; \emph{emptiness} folds through \texttt{simd\_any}.
\end{description}

\begin{figure}[t]
\centering
\begin{tikzpicture}[font=\scriptsize, node distance=3mm,
  cell/.style={draw, minimum size=3.1mm, inner sep=0pt, fill=cGrey!12},
  own/.style={cell, fill=cA!35, draw=cA},
  own2/.style={cell, fill=cC!35, draw=cC},
  lab/.style={font=\small\bfseries},
  mem/.style={draw, rounded corners=1.5pt, font=\scriptsize, minimum height=5mm}]

% R1
\node[lab] at (1.1,2.35) {R1 private};
\foreach \i in {0,...,5} \foreach \j in {0,...,5}
  \node[own] at (0.35*\j+0.2, 1.9-0.35*\i) {};
\node[mem, fill=cB!12, draw=cB, below=1mm of {(1.1,0.0)}, anchor=north] (m1)
  {\textbf{device memory} (spilled)};
\node[below=0.5mm of m1, align=center] {$\mathit{dim}^2$ per \emph{thread}\\2304\,B: no register fits};

% R2
\node[lab] at (4.5,2.35) {R2 threadgroup};
\foreach \i in {0,...,5} \foreach \j in {0,...,5}
  \node[own] at (0.35*\j+3.6, 1.9-0.35*\i) {};
\node[mem, fill=cC!12, draw=cC, below=1mm of {(4.5,0.0)}, anchor=north] (m2)
  {\textbf{on-chip} 32\,KB budget};
\node[below=0.5mm of m2, align=center] {$32\times576$\,B $=18$\,KB\\$\Rightarrow$ 1 resident group};

% R3
\node[lab] at (8.1,2.35) {R3 row-distributed};
\foreach \i in {0,...,5} \foreach \j in {0,...,5}
  \node[cell] at (0.35*\j+7.2, 1.9-0.35*\i) {};
\foreach \j in {0,...,5} \node[own2] at (0.35*\j+7.2, 1.9-0.35*2) {};
\node[left=1mm of {(7.15,1.9-0.7)}, cC, font=\scriptsize] {lane $i$};
\node[mem, fill=cA!12, draw=cA, below=1mm of {(8.1,0.0)}, anchor=north] (m3)
  {\textbf{registers} + \texttt{simd\_shuffle}};
\node[below=0.5mm of m3, align=center] {$\mathit{dim}$ per \emph{lane}\\$24\times$ smaller};
\end{tikzpicture}
\caption{Three DBM representations. R3 solves the footprint problem outright, and still loses
--- for the reason given in Sect.~\ref{sec:cooperation}.}
\label{fig:reps}
\end{figure}

\subsection{Adaptive CPU/GPU merge routing}
\label{sec:adaptive}

Per-discrete-state merging (\aLU{} subsumption plus convex union) is a second kernel: one
bucket per thread. Because bucket shape varies enormously between models, we route each BFS
level to whichever device is predicted faster. The router calibrates the throughput of each
path on early levels --- measuring $\mathit{work}/\mathit{ms}$ with
$\mathit{work} = \sum_b |b|^2$, the natural cost model for pairwise antichain folding ---
then routes each later level to the higher-throughput device, updating its estimate with an
exponential moving average (Fig.~\ref{fig:adaptive}).

The router is useful in its own right, but its scientific value here is as an
\emph{instrument}: given a free choice on every level, its selections are an unbiased
verdict on which device is actually faster.

\begin{figure}[t]
\centering
\begin{tikzpicture}[font=\small,
  b/.style={draw, rounded corners=2pt, minimum height=7.5mm, align=center, font=\scriptsize},
  ar/.style={-{Latex[length=1.8mm]}, thick}]
\node[b, fill=cGrey!12] (lvl) {level $k$\\buckets};
\node[b, right=7mm of lvl, fill=cD!12, draw=cD] (cal) {$r_{\mathit{cpu}},r_{\mathit{gpu}}$\\known?};
\node[b, above right=3mm and 8mm of cal, fill=cB!12, draw=cB] (probe) {probe the\\uncalibrated path};
\node[b, below right=3mm and 8mm of cal, fill=cC!12, draw=cC] (route) {route to
$\arg\max r$};
\node[b, right=8mm of route, fill=cGrey!12] (upd) {update $r$\\(EMA)};
\draw[ar] (lvl) -- (cal);
\draw[ar] (cal) -- node[above left, font=\scriptsize] {no} (probe);
\draw[ar] (cal) -- node[below left, font=\scriptsize] {yes} (route);
\draw[ar] (route) -- (upd);
\draw[ar] (probe.east) -- ++(6mm,0) |- (upd.north);
\node[right=2mm of upd, font=\scriptsize, align=left, cB]
  {measured outcome:\\\textbf{122/124 levels $\to$ CPU}};
\end{tikzpicture}
\caption{Adaptive per-level merge routing. Used as a measurement instrument, it selects the
GPU only for the single calibration probe on each model.}
\label{fig:adaptive}
\end{figure}
